\SourceLeaf{464}
\par\Needspace{8\baselineskip}\bigskip
\pdfbookmark[3]{Exposition, 15}{pallas-II-article15}
\begin{center}15.\end{center}
\noindent\hspace*{\parindent}Les recherches précédentes sont indépendantes de la nature des forces perturbatrices. Supposons maintenant, que celles-ci sont produites par l’action qu’un autre corps exerce tant sur la planète que sur le soleil. Soit \(\mu m\) la masse de ce corps; \(x',y',z'\) ses coordonnées par rapport aux trois plans fondamentaux; \(r'=\BareRoot{(x'x'+y'y'+z'z')}\) sa distance au soleil; \(\rho=\BareRoot{((x'-x)^2+(y-y')^2+(z-z')^2)}\) sa distance à la planète troublée. On aura donc d’après les principes connues
\[
\begin{aligned}
\xi&=\dfrac{\mu mx'}{r'^3}-\dfrac{\mu m(x'-x)}{\rho^3}\\
\eta&=\dfrac{\mu my'}{r'^3}-\dfrac{\mu m(y'-y)}{\rho^3}\\
\zeta&=\dfrac{\mu mz'}{r'^3}-\dfrac{\mu m(z'-z)}{\rho^3}\mathpunct{\cdot}
\end{aligned}
\]
\par En désignant par \(L'\) le lieu héliocentrique de la planète perturbante, nous aurons
\[
\begin{aligned}
x'&=r'\cos XL'\\
y'&=r'\cos YL'\\
z'&=r'\cos ZL';
\end{aligned}
\]
ainsi les formules de l’article 9 donnent (en faisant attention que \(\Lambda L=ML=90^0\)):
\[
\begin{aligned}
T&=\mu\left(\dfrac{1}{\rho^3}-\dfrac{1}{r'^3}\right)r'\cos LL'-\dfrac{\mu r}{\rho^3}\\
V&=\mu\left(\dfrac{1}{\rho^3}-\dfrac{1}{r'^3}\right)r'\cos\Lambda L'\\
W&=\mu\left(\dfrac{1}{\rho^3}-\dfrac{1}{r'^3}\right)r'\cos ML'.
\end{aligned}
\]
Soit \(\beta'\) l’inclinaison du rayon vecteur \(r'\) au plan de l’orbite osculatrice de la planète troublée (prise avec le signe \(+\) du côté où est le pôle \(M\)), et que sa projection sur ce plan fasse, avec la ligne du noeud ascendant du même plan sur le plan des coordonnées \(x,y\), l’angle \(v'-\AscendingNode\), pris dans le sens du mouvement de la planète troublée, de sorte que \(\beta',v'\) puissent être considérées comme latitude et longitude héliocentrique du corps perturbant relativement \SourceLeaf{465}au plan de l’orbite osculatrice de la planète troublée. Il sera facile de voir qu’on a
\[
\begin{aligned}
\cos LL'&=\cos\beta'\cos(v'-v)\\
\cos\Lambda L'&=\cos\beta'\sin(v'-v)\\
\cos ML'&=\sin\beta'.
\end{aligned}
\]
\par Enfin pour réunir ici tout ce dont on a besoin pour le calcul, supposons l’inclinaison du plan, dans lequel se meut le corps perturbant, au plan des coordonnées \(x,y\) égale à \(i'\), et que \(\AscendingNode'\) soit la longitude du noeud ascendant comptée de l’axe des \(x\); de plus soit \(w'-\AscendingNode'\) la distance du lieu héliocentrique \(L'\) à ce même noeud, ou \(w'\) ce qu’on appelle sa longitude vraie dans l’orbite. Qu’on détermine les quantités \(N,\Delta,J\) au moyen des équations suivantes:
\[
\begin{aligned}
\sin\tfrac12J\mathbin{.}\sin\tfrac12(\Delta+N)&=\sin\tfrac12(i'+i)\mathbin{.}\sin\tfrac12(\AscendingNode'-\AscendingNode)\\
\sin\tfrac12J\mathbin{.}\cos\tfrac12(\Delta+N)&=\sin\tfrac12(i'-i)\mathbin{.}\cos\tfrac12(\AscendingNode'-\AscendingNode)\\
\cos\tfrac12J\mathbin{.}\sin\tfrac12(\Delta-N)&=\cos\tfrac12(i'+i)\mathbin{.}\sin\tfrac12(\AscendingNode'-\AscendingNode)\\
\cos\tfrac12J\mathbin{.}\cos\tfrac12(\Delta-N)&=\cos\tfrac12(i'-i)\mathbin{.}\cos\tfrac12(\AscendingNode'-\AscendingNode).
\end{aligned}
\]
\par Cela fait on aura
\[
\begin{aligned}
\operatorname{tang}(v'-\AscendingNode-\Delta)&=\cos J\mathbin{.}\operatorname{tang}(w'-\AscendingNode'-N)\\
\operatorname{tang}\beta'&=\operatorname{tang}J\mathbin{.}\sin(v'-\AscendingNode-\Delta)\\
\sin\beta'&=\sin J\mathbin{.}\sin(w'-\AscendingNode'-N)\\
\rho&=\BareRoot{(rr+r'r'-2rr'\cos\beta'\cos(v'-v))}.
\end{aligned}
\]
